\section{Introduction}
\label{sec:Introduction}
Scientific workflow engines reserve memory before tasks execute, despite
variation in their inputs and behavior. Insufficient reservations risk
out-of-memory (OOM) failures and retries; excessive reservations occupy
capacity that could serve other tasks. Prediction must balance these costs.

A point predictor commonly estimates a conditional mean or median, which
describes typical demand but does not by itself specify an upper-tail
reservation. Quantile regression instead targets chosen percentiles of
the conditional memory distribution~\cite{meinshausen2006quantile}.
Each quantile estimate is a single value, but several quantiles describe
central demand and the upper tail. An upper quantile directly targets
underallocation risk; realized coverage still depends on model quality,
calibration, and input changes.

Existing online allocators address this risk through different mechanisms.
Sizey selects or combines regression predictions and adds an adaptive
offset~\cite{bader2024sizey}. Ponder switches between historical bounds,
user defaults, and regression that penalizes underprediction more strongly,
with additional safeguards~\cite{lehmann2024ponder}. Their described sizing
methods use input size and execution history, but the amount of data a
pending task will actually consume is unknown before execution. Staged size need not
describe actual accesses: a tool may touch only part of a mapped index or
stream many bytes while retaining little memory. Consumption supplies
context beyond input size, rather than measuring peak memory itself.

Consumption measurement traces bytes read and estimates memory-mapped
accesses from file-backed page faults attributed to each task and its
child processes. This auditing incurs tracing and processing costs.
Selective auditing uses a gating process to select tasks whose consumption
features are expected to improve memory predictions for similar subsequent
tasks, while respecting an audit budget.

Our prior WORKS paper introduced Consumption-Aware Memory Prediction
(\toolname)~\cite{camp2026works} to combine these mechanisms. CAMP uses
static task features, historical peak memory, and consumption feedback
from task-attributed eBPF auditing. An auxiliary model predicts consumption
before launch; hierarchical quantile models and residual calibration then
produce a memory request. Selective auditing supplies further feedback.
CAMP showed that combining quantile modeling with data-consumption
features can improve memory prediction accuracy over Sizey~\cite{camp2026works}.
Its architecture, ablations, and selective-auditing results remain prior work.

This paper expands analysis of those component measurements and studies
live deployment on repeated and previously unseen Bowtie2 inputs. Gains
on familiar inputs need not persist when inputs change. Does an extra
warm pass improve allocation on unseen inputs?
Warm-Transfer processes the new inputs after that pass; Direct-Transfer
starts from the same cold checkpoint without it. Both branches continue
learning online. These live runs audit every task; selective auditing is
evaluated separately using retained measurements.

Each stage executes 6,000 tasks. Repeated inputs reduce requests by
75.34\%. On new inputs, Warm-Transfer uses less unused memory-time
(0.12805 versus 0.22109~GiB-hours), but has lower request coverage
(94.15\% versus 95.78\%). We therefore examine requested capacity,
unused memory-time, and first-attempt request margins together.

Our evaluation comprises:
\begin{itemize}
\item \emph{Measurement and features:} task filtering and workflow
integration, read-only versus read-plus-mmap measurements, and four
consumption-feature variants;
\item \emph{Downstream allocation and auditing cost:} expanded analysis
of retained selective-versus-random auditing measurements and separate
auditing-cost estimates; and
\item \emph{Online deployment and adaptation:} repeated-input learning
and new-input transfer with and without an additional warm pass.
\end{itemize}
Sections~\ref{sec:problem}--\ref{sec:design} present the problem and design;
Sections~\ref{sec:experiments}--\ref{sec:results}, the evaluation; and
Sections~\ref{sec:related}--\ref{sec:conclusion}, related work and conclusions.
